\documentclass[10pt,a4paper]{article}
\usepackage{fontspec}
\setmainfont{Liberation Sans}
\usepackage{polyglossia}
\setdefaultlanguage{polish}
\usepackage[a4paper,margin=19mm,headheight=15pt]{geometry}
\usepackage{xcolor,tabularx,booktabs,array,enumitem,fancyhdr,titlesec,hyperref}
\definecolor{navy}{HTML}{263650}
\definecolor{blue}{HTML}{5475A5}
\definecolor{pale}{HTML}{EFF3F8}
\definecolor{muted}{HTML}{5D6878}
\hypersetup{colorlinks=true,urlcolor=blue}
\pagestyle{fancy}
\fancyhf{}
\fancyhead[L]{\small\color{muted}Tutorland}
\fancyhead[R]{\small\color{muted}Mapa stron}
\fancyfoot[C]{\small\color{muted}\thepage}
\renewcommand{\headrulewidth}{0.3pt}
\setlength{\parindent}{0pt}
\setlength{\parskip}{5pt}
\setlength{\emergencystretch}{3em}
\sloppy
\titleformat{\section}{\Large\bfseries\color{navy}}{\thesection.}{0.55em}{}
\newcolumntype{Y}{>{\raggedright\arraybackslash}X}
\begin{document}
\begin{titlepage}
\thispagestyle{empty}
\vspace*{32mm}
{\large\color{blue}\textbf{DOCELOWA STRUKTURA PRODUKTU}}
\vspace{12mm}
{\Huge\bfseries\color{navy}Tutorland\par}
\vspace{4mm}
{\LARGE Wszystkie potrzebne strony i widoki\par}
\vspace{13mm}
\color{blue}\rule{\textwidth}{2pt}
\vspace{9mm}
{\large Mapa publicznych ekranów oraz wspólnych widoków prywatnych w finalnej wersji projektu.}
\vfill
\colorbox{pale}{\parbox{\dimexpr\textwidth-2\fboxsep\relax}{\textbf{Jak czytać dokument}\\
Obecny prototyp zawiera katalog, profil korepetytora, filtry i kalendarz rezerwacji. Pozostałe ekrany poniżej to docelowy zakres do zaprojektowania i wdrożenia; część pozycji może być widokiem lub stanem jednej strony, a nie osobną trasą URL.}}\par
\vspace{12mm}
{\small\color{muted}Struktura oparta na zakresie produktu Tutorland: uczniowie i korepetytorzy, rezerwacje 40-minutowe, Zoom, spotkania i opinie.}
\end{titlepage}

\section{Strony publiczne}
\small
\begin{tabularx}{\textwidth}{>{\bfseries}p{8mm}p{41mm}p{30mm}Y}
\toprule
Nr & Strona / przykładowa trasa & Dostęp & Zawartość i działanie\\
\midrule
1 & Strona główna i katalog\newline \url{/} lub \url{/korepetytorzy} & Każdy & Wprowadzenie do usługi, lista ofert, wyszukiwanie i filtry przedmiotów, poziomu, ceny, doświadczenia i dostępności; stan pustych wyników. To główny punkt wejścia.\\
2 & Profil korepetytora\newline \url{/korepetytorzy/:slug} & Każdy & Opis, przedmioty, poziomy, cena, doświadczenie, oceny i opinie; kalendarz dostępności i początek rezerwacji. Przy zapisie użytkownik niezalogowany przechodzi do logowania, a po nim wraca do wyboru terminu.\\
\bottomrule
\end{tabularx}
\normalsize

\section{Konto i wspólne widoki prywatne}
\small
\begin{tabularx}{\textwidth}{>{\bfseries}p{8mm}p{41mm}p{30mm}Y}
\toprule
Nr & Strona / przykładowa trasa & Dostęp & Zawartość i działanie\\
\midrule
3 & Rejestracja\newline \url{/rejestracja} & Gość & Utworzenie konta i wybór roli: uczeń albo korepetytor; walidacja danych, zgoda na wymagane warunki i czytelny wynik.\\
4 & Logowanie\newline \url{/logowanie} & Gość & Logowanie i bezpieczny powrót do strony lub rezerwacji, z której rozpoczęto logowanie.\\
5 & Ustawienia konta\newline \url{/konto/ustawienia} & Zalogowany & Dane profilu, zmiana hasła i ustawień; użytkownik może edytować wyłącznie własne dane.\\
6 & Spotkania\newline \url{/spotkania} & Uczeń / korepetytor & Lista przyszłych i zakończonych lekcji; data, godzina, przedmiot, status, dane drugiej strony, link Zoom oraz powiadomienia w aplikacji. Zawiera stany pustej listy i błędu.\\
7 & Formularz opinii\newline \url{/spotkania/:id/opinia} lub stan na stronie spotkania & Uczeń po lekcji & Ocena 1–5 i komentarz do własnej zakończonej rezerwacji. Jedna opinia na rezerwację; formularz nieaktywny przed planowanym końcem lekcji.\\
\bottomrule
\end{tabularx}
\normalsize
\end{document}
